% =============================================================================
%  FPGA Accelerated Aging Environment — Student Onboarding Guide
%  Laboratory of Electronic Systems and Circuits (LESC)
%  Federal University of Ceará
% =============================================================================
\documentclass[12pt, a4paper]{report}

\usepackage[a4paper, margin=2.5cm, headheight=15pt]{geometry}
\usepackage{fontenc}
\usepackage{inputenc}
\usepackage{lmodern}
\usepackage{microtype}
\usepackage{hyperref}
\usepackage{xcolor}
\usepackage{listings}
\usepackage{booktabs}
\usepackage{tabularx}
\usepackage{amsmath}
\usepackage{amssymb}
\usepackage{graphicx}
\usepackage{float}
\usepackage{fancyhdr}
\usepackage{titlesec}
\usepackage{enumitem}
\usepackage{parskip}
\usepackage{mdframed}
\usepackage{tcolorbox}
\usepackage{array}
\usepackage{longtable}
\usepackage[T1]{fontenc}

% ---------------------------------------------------------------------------
% Colours
% ---------------------------------------------------------------------------
\definecolor{codeBackground}{RGB}{40, 42, 54}
\definecolor{codeForeground}{RGB}{248, 248, 242}
\definecolor{keywordColor}{RGB}{255, 121, 198}
\definecolor{commentColor}{RGB}{98, 114, 164}
\definecolor{stringColor}{RGB}{241, 250, 140}
\definecolor{warningBox}{RGB}{255, 243, 205}
\definecolor{infoBox}{RGB}{209, 236, 241}
\definecolor{successBox}{RGB}{212, 237, 218}
\definecolor{headerGray}{RGB}{68, 71, 90}

% ---------------------------------------------------------------------------
% Code listings
% ---------------------------------------------------------------------------
\lstdefinestyle{shell}{
  backgroundcolor=\color{codeBackground},
  basicstyle=\ttfamily\footnotesize\color{codeForeground},
  keywordstyle=\color{keywordColor}\bfseries,
  commentstyle=\color{commentColor},
  stringstyle=\color{stringColor},
  breaklines=true,
  breakatwhitespace=true,
  frame=single,
  framerule=0.5pt,
  rulecolor=\color{headerGray},
  xleftmargin=6pt,
  xrightmargin=6pt,
  showstringspaces=false,
  tabsize=2,
}

\lstdefinestyle{verilog}{
  backgroundcolor=\color{codeBackground},
  basicstyle=\ttfamily\footnotesize\color{codeForeground},
  keywordstyle=\color{keywordColor}\bfseries,
  commentstyle=\color{commentColor},
  stringstyle=\color{stringColor},
  language=Verilog,
  breaklines=true,
  frame=single,
  framerule=0.5pt,
  rulecolor=\color{headerGray},
  xleftmargin=6pt,
  xrightmargin=6pt,
  showstringspaces=false,
  tabsize=2,
}

\lstdefinestyle{python}{
  backgroundcolor=\color{codeBackground},
  basicstyle=\ttfamily\footnotesize\color{codeForeground},
  keywordstyle=\color{keywordColor}\bfseries,
  commentstyle=\color{commentColor},
  stringstyle=\color{stringColor},
  language=Python,
  breaklines=true,
  frame=single,
  framerule=0.5pt,
  rulecolor=\color{headerGray},
  xleftmargin=6pt,
  xrightmargin=6pt,
  showstringspaces=false,
  tabsize=4,
}

% ---------------------------------------------------------------------------
% Coloured boxes
% ---------------------------------------------------------------------------
\tcbuselibrary{skins, breakable}

\newtcolorbox{warningbox}[1][]{
  colback=warningBox, colframe=orange!80!black,
  fonttitle=\bfseries, title=Warning, #1
}

\newtcolorbox{infobox}[1][]{
  colback=infoBox, colframe=blue!50!black,
  fonttitle=\bfseries, title=Note, #1
}

\newtcolorbox{tipbox}[1][]{
  colback=successBox, colframe=green!50!black,
  fonttitle=\bfseries, title=Tip, #1
}

% ---------------------------------------------------------------------------
% Page style
% ---------------------------------------------------------------------------
\pagestyle{fancy}
\fancyhf{}
\fancyhead[L]{\small FPGA Aging Environment — LESC / UFC}
\fancyhead[R]{\small \leftmark}
\fancyfoot[C]{\small \thepage}
\renewcommand{\headrulewidth}{0.4pt}

% ---------------------------------------------------------------------------
% Section formatting
% ---------------------------------------------------------------------------
\titleformat{\chapter}[block]
  {\normalfont\Large\bfseries\color{headerGray}}
  {\thechapter.}{10pt}{}
\titlespacing*{\chapter}{0pt}{20pt}{12pt}

% ---------------------------------------------------------------------------
% Hyperref setup
% ---------------------------------------------------------------------------
\hypersetup{
  colorlinks=true,
  linkcolor=blue!60!black,
  urlcolor=blue!70!black,
  citecolor=green!50!black,
  pdftitle={FPGA Aging Environment — Onboarding Guide},
  pdfauthor={LESC / UFC},
}

% =============================================================================
\begin{document}

% ---------------------------------------------------------------------------
% Title page
% ---------------------------------------------------------------------------
\begin{titlepage}
  \centering
  \vspace*{2cm}
  {\Large\bfseries Federal University of Ceará\\[6pt]
   Department of Electrical Engineering\\[6pt]
   Laboratory of Electronic Systems and Circuits (LESC)\\[2cm]}
  \rule{\linewidth}{1pt}\\[0.5cm]
  {\Huge\bfseries FPGA Accelerated Aging Environment\\[0.4cm]}
  {\Large Student Onboarding Guide\\[0.3cm]}
  \rule{\linewidth}{1pt}\\[2cm]
  {\large A complete reference for understanding, setting up, and operating\\
   the hardware-in-the-loop burn-in experiment platform.\\[3cm]}
  {\normalsize Revision: May 2026\\[1cm]}
  \vfill
\end{titlepage}

\tableofcontents
\newpage

% =============================================================================
\chapter{Introduction}
% =============================================================================

This document describes the \textbf{FPGA Accelerated Aging Environment} developed at
LESC (Laboratory of Electronic Systems and Circuits, Federal University of Ceará).
Its purpose is to give a new student or contributor everything needed to understand
the scientific motivation, operate the equipment, and extend the software or RTL.

\section{What the project does}

Semiconductor transistors degrade over time when subjected to electrical stress.
The key mechanisms --- Negative Bias Temperature Instability (NBTI) and Hot Carrier
Injection (HCI) --- cause the threshold voltage $V_{th}$ of MOSFET devices to drift,
increasing propagation delays in logic gates. If degradation is severe enough, a once-working
digital circuit will begin to fail timing constraints and produce erroneous outputs.

Studying these phenomena experimentally requires either very long stress times at
nominal conditions, or \emph{accelerated aging}: raising the temperature and/or supply
voltage to stress the device much faster. This project implements a complete hardware-in-the-loop
environment for accelerated aging of commercial FPGAs.

The \emph{device under test} (DUT) is a Digilent Nexys4 DDR board populated with a
Xilinx Artix-7 FPGA (xc7a100tcsg324-1). It is placed inside a temperature-controlled
oven, powered by a programmable benchtop supply, and continuously monitored by a
custom RTL bitstream that measures timing degradation in real time using a metastability-based
phase sensor. All measurements are streamed over UART to a host PC where a Python
application logs timestamped CSV files and displays live plots.

\section{Who this document is for}

This guide assumes the reader has:
\begin{itemize}
  \item Familiarity with digital design concepts (clock domains, setup/hold time, FPGAs).
  \item Basic SystemVerilog or Verilog reading ability.
  \item Experience with Python and comfort at the Linux command line.
  \item Some exposure to control systems (PID loops, transfer functions).
\end{itemize}

No prior knowledge of the project is required.

\section{Document scope}

Chapter~\ref{ch:background} covers the device-physics background --- NBTI, HCI, and the
reaction-diffusion model --- at a graduate level. Chapter~\ref{ch:system} gives a
system-level overview. Chapters~\ref{ch:rtl} and~\ref{ch:software} describe the FPGA
RTL and the Python software in depth. Chapter~\ref{ch:protocol} documents all
communication protocols. Chapter~\ref{ch:control} explains the control loops.
Chapter~\ref{ch:setup} provides step-by-step workstation setup instructions.
Chapters~\ref{ch:running} and~\ref{ch:analysis} cover running an experiment and
analyzing the results.

% =============================================================================
\chapter{Scientific Background}
\label{ch:background}
% =============================================================================

\section{Why FPGAs age}

A modern FPGA is built from hundreds of thousands of CMOS transistors. Each transistor
consists of a gate electrode separated from the silicon channel by a thin gate oxide.
The electric field across this oxide, combined with the kinetic energy of charge carriers,
gradually damages the oxide and the Si/SiO\textsubscript{2} interface. The primary
mechanisms that affect digital timing are NBTI and HCI.

\section{Negative Bias Temperature Instability (NBTI)}

NBTI is the dominant aging mechanism in modern sub-micron CMOS. It predominantly
affects \textbf{PMOS} transistors biased in inversion (gate-to-source voltage $V_{GS} < 0$),
which is precisely the stressed condition during logic `1' propagation through a PMOS pull-up.

\subsection{Microscopic mechanism: interface trap generation}

The Si/SiO\textsubscript{2} interface is chemically passivated by hydrogen. Silicon
dangling bonds --- known as $P_b$ centres --- are tied up as Si-H bonds at the interface.
Under a strong negative gate bias and elevated temperature, these bonds dissociate:

\begin{equation}
  \text{Si} \text{---} \text{H} \;\xrightarrow{\,V_{GS},\,T\,}\; \text{Si}\bullet \;+\; \text{H}
\end{equation}

The resulting silicon dangling bond becomes a positively charged interface trap that
captures a hole and shifts the PMOS threshold voltage:

\begin{equation}
  \Delta V_{th} = \frac{q \, N_{it}}{C_{ox}}
\end{equation}

where $N_{it}$ is the interface trap density, $q$ is the electron charge, and $C_{ox}$
is the oxide capacitance per unit area. In addition, some hydrogen and holes are captured
as \emph{oxide traps} in the bulk SiO\textsubscript{2}, contributing a further positive
charge that also increases $|V_{th}|$.

An increased $|V_{th}|$ weakens the on-current of the PMOS ($I_{D} \propto V_{GS} - V_{th}$),
reducing its drive strength and increasing the propagation delay of every gate where
a PMOS pull-up is on the critical path.

\subsection{Reaction-Diffusion model}

The temporal evolution of $\Delta V_{th}$ is described by the Reaction-Diffusion (R-D) model.
During the \textbf{stress phase}, interface traps are generated at the interface (reaction)
and hydrogen diffuses away into the oxide (diffusion), preventing back-passivation. The
rate-limiting step is the diffusion of atomic or molecular hydrogen. Solving the
coupled reaction-diffusion equations yields the well-known power-law time dependence:

\begin{equation}
  \Delta V_{th}(t) = A \left(\frac{V_{GS,stress}}{t_{ox}}\right)^{m} \exp\!\left(\frac{-E_a}{k_B T}\right) \cdot t^{n}
\end{equation}

where:
\begin{itemize}
  \item $n \approx 0.1$--$0.25$ is the time exponent (determined by whether H or H\textsubscript{2} diffuses);
  \item $E_a \approx 0.1$--$0.2\,\text{eV}$ is the activation energy for the diffusion process;
  \item $k_B$ is Boltzmann's constant and $T$ is absolute temperature;
  \item $m$ is the field-acceleration exponent;
  \item $A$ is a technology-dependent pre-factor.
\end{itemize}

\subsection{Recovery}

A key (and sometimes inconvenient) feature of NBTI is that degradation is
\emph{partially reversible}. When the gate bias is removed, residual hydrogen
at the interface back-passivates some of the dangling bonds:
$\text{Si}\bullet \;+\; \text{H} \;\rightarrow\; \text{Si-H}$.
This \textbf{recovery} can mask the true degradation if measurements are taken
after a rest period. In the long run, however, the degradation is not fully
recovered --- the component that does not recover (often called the permanent
component) accumulates monotonically and represents the true aging damage.

In this experiment the DUT is stressed \emph{continuously} (no recovery periods),
so the measured $\Delta V_{th}$ is representative of the full stress-only trajectory.

\subsection{Temperature acceleration}

The Arrhenius relationship gives the ratio of experiment lifetimes at two temperatures:

\begin{equation}
  \text{AF} = \frac{\tau_{use}}{\tau_{stress}} =
  \exp\!\left[\frac{E_a}{k_B}\left(\frac{1}{T_{use}} - \frac{1}{T_{stress}}\right)\right]
\end{equation}

For $E_a = 0.15\,\text{eV}$, stressing at $125\,^\circ\text{C}$ instead of $25\,^\circ\text{C}$
gives an acceleration factor of roughly $8\times$. For higher activation energies the
acceleration is even more pronounced. This is why the experiment uses an oven: a DUT
at $125\,^\circ\text{C}$ sees months of aging compressed into days.

\section{Hot Carrier Injection (HCI)}

HCI predominantly affects \textbf{NMOS} transistors operating at high drain-to-source
voltage $V_{DS}$. Near the drain, the lateral electric field $\mathcal{E} \approx V_{DS}/L$
is highest. Electrons accelerating towards the drain can acquire kinetic energies
well above the thermal average (``hot'' carriers). Some fraction of these hot carriers
have sufficient energy to surmount the Si/SiO\textsubscript{2} barrier ($\approx 3.1\,\text{eV}$)
and be injected into the gate oxide, where they are trapped or generate interface states.

The resulting threshold voltage shift and transconductance degradation slow NMOS
pull-down paths, directly increasing fall times in logic gates. HCI is worst at
intermediate $V_{GS}$ values that maximise the substrate current (maximum multiplication
condition), and it accelerates strongly with supply voltage.

In this experiment, the supply voltage $V_{CCINT}$ is elevated above the nominal
$1.0\,\text{V}$ to accelerate both NBTI and HCI simultaneously. The trade-off is that
the acceleration factor must be chosen carefully: too high a voltage can trigger
catastrophic dielectric breakdown rather than controlled, gradual aging.

\section{How degradation manifests in digital timing}

From a circuit perspective, NBTI and HCI both increase propagation delay $t_{pd}$
along paths that traverse affected transistors. In a synchronous design, the critical
constraint is the \emph{setup time} relation:

\begin{equation}
  t_{pd} \;\leq\; T_{clk} - t_{setup} - t_{skew}
\end{equation}

The \emph{timing slack} is the margin available:

\begin{equation}
  \text{slack} = T_{clk} - t_{setup} - t_{skew} - t_{pd}
\end{equation}

As $t_{pd}$ increases due to aging, slack decreases. When slack reaches zero, the
flip-flop at the path endpoint can no longer reliably capture the correct data value,
and functional errors begin to occur.

This project measures timing degradation by directly quantifying this slack using an
on-chip phase-sensitive sensor, described in detail in Chapter~\ref{ch:rtl}.

\section{Why FPGAs are useful for aging experiments}

Commercial FPGAs offer several advantages as aging test vehicles:

\begin{enumerate}
  \item \textbf{Fixed routing}: the FPGA's routing fabric is programmed once (fixed PnR), so the
        same physical transistors are stressed repeatedly. Unlike an ASIC, there is no process-to-process
        variation between samples.
  \item \textbf{Reprogrammability}: the bitstream can be changed to instrument different paths or
        add new sensors without re-fabricating silicon.
  \item \textbf{Built-in XADC}: Xilinx FPGAs include an on-chip Analogue-to-Digital Converter (XADC)
        that continuously measures die temperature and supply voltages without external probes.
  \item \textbf{MMCM phase control}: the Mixed-Mode Clock Manager (MMCM) supports dynamic phase
        shift via the DRP interface, enabling sub-nanosecond timing manipulation at runtime.
\end{enumerate}

% =============================================================================
\chapter{System Overview}
\label{ch:system}
% =============================================================================

\section{Architecture block diagram}

The complete system is shown below. Three separate subsystems communicate with the host PC:

\begin{lstlisting}[style=shell, caption={System block diagram (ASCII)}]
                     +-----------------+
                     |   Host PC       |
                     |  (App_Nexys)    |
                     +---+---+---+-----+
                         |   |   |
      Serial 9600 baud   |   |   |  USB-TMC (VISA)
  .-----------------------'   |   '------------------------.
  |                           |  Serial 115200 baud        |
  v                           v                            v
+---------------+    +-------------------+      +--------------------+
|  Nexys4 DDR   |    |  Arduino UNO/Mega |      |  ITECH IT6502D     |
|  (FPGA DUT)   |    |  (Oven PID)       |      |  (Programmable PSU)|
|               |    |                   |      |                    |
|  Artix-7      |    |  NTC thermistor   |      |  Controls VCCINT   |
|  xc7a100t     |    |  SSR output       |      |  of the FPGA       |
|               |    |                   |      |                    |
+--+------------+    +---------+---------+      +---------+----------+
   |                           |                          |
   |  Inside oven              |                          |
   |  (elevated temp)          v                          |
   |                  +--------+--------+                 |
   |                  |  Electric Oven  |                 |
   |                  |  + SSR          |                 |
   |                  |  + NTC probe    |                 |
   |                  +-----------------+                 |
   |                                                      |
   '---------- FPGA board is inside the oven -------------'
                (Arduino/PSU outside, only UART cables enter)
\end{lstlisting}

\begin{infobox}
The Arduino and PSU are optional. The app can run with only the FPGA connected
(no temperature control, no programmable supply). However, meaningful aging experiments
require all three subsystems.
\end{infobox}

\section{Hardware bill of materials}

\begin{longtable}{@{}lll@{}}
  \toprule
  \textbf{Component} & \textbf{Role} & \textbf{Interface} \\
  \midrule
  \endhead
  Digilent Nexys4 DDR & DUT (Artix-7 FPGA) & USB-UART (FTDI) \\
  Electric oven / incubator & Thermal stress chamber & --- \\
  Solid-State Relay (SSR) & Oven power switching & Digital GPIO (Arduino) \\
  NTC thermistor (10~k$\Omega$) & Oven temperature sensing & Analogue input (Arduino) \\
  Arduino UNO or Mega & PID controller & USB-UART, 115200 baud \\
  ITECH IT6502D & Programmable DC supply & USB-TMC (VISA) \\
  Host PC (Linux) & Data acquisition \& control & USB hub \\
  \bottomrule
\end{longtable}

\section{Experiment overview}

A typical aging experiment proceeds as follows:

\begin{enumerate}
  \item The FPGA board is placed inside the oven. USB/UART cables pass through
        a small opening in the oven door.
  \item The IT6502D PSU is connected to the FPGA's external power header to supply
        an elevated $V_{CCINT}$.
  \item The host PC is powered on; all USB devices enumerate.
  \item The operator launches the \texttt{App\_Nexys} application and configures the
        serial ports in the Setup Dialog.
  \item The operator sets the oven temperature setpoint (e.g., $125\,^\circ\text{C}$)
        and the target $V_{CCINT}$ (e.g., $1.1\,\text{V}$).
  \item The operator clicks \textbf{Start Test}. The application enters autonomous mode:
        \begin{itemize}
          \item The Arduino PID loop ramps the oven to the setpoint.
          \item The PSU delivers the commanded voltage; the closed-loop trim adjusts it
                to match the XADC-measured $V_{CCINT}$.
          \item Every second, the PC sends a trigger byte (\texttt{0x54}) to the FPGA.
          \item The FPGA responds with a 15-byte packet: die temperature, slack counter,
                $V_{CCINT}$, failure flag, and canary metrics.
          \item All data is appended to a timestamped CSV file.
        \end{itemize}
  \item The experiment runs unattended for one to seven days.
  \item The operator clicks \textbf{Stop Test} and retrieves the CSV for analysis.
\end{enumerate}

% =============================================================================
\chapter{FPGA RTL Architecture}
\label{ch:rtl}
% =============================================================================

\section{Target device}

\begin{tabular}{ll}
  Board: & Digilent Nexys4 DDR \\
  Device: & Xilinx Artix-7 xc7a100tcsg324-1 \\
  Top module: & \texttt{nexys4\_aging\_top} \\
  Source language: & SystemVerilog (pure RTL, no block design) \\
  Toolchain: & Vivado 2025.2+ \\
\end{tabular}

\section{Module hierarchy}

\begin{lstlisting}[style=shell, caption={RTL module hierarchy}]
nexys4_aging_top                    (src/rtl/top/)
  clk_wiz_0                         MMCM -- dynamic phase shift (src/ip/)
  XADC                              on-chip ADC (primitive)
  temp_catcher                      XADC DRP monitor (aging_sensor/)
  modern_sensible                   3-FF metastability sensor (aging_sensor/)
  [2-FF synchronizer]               alarm CDC from clk_en to clk_sys
  adder_canary                      dual-adder sensor (aging_sensor/)
    u_sensor  (ripple_adder)          toggle-driven, timing measurement
    u_canary  (ripple_adder)          counter-driven, functional canary
  controller_controller             phase-sweep FSM (aging_sensor/)
  failure_holder                    sticky failure latch (aging_sensor/)
  uart_rx                           PC -> FPGA trigger receiver (uart/)
  sensor_stream                     FPGA -> PC packet serialiser (uart/)
  uart_tx                           UART transmitter (uart/)
  BINtoBCD (x2)                     for 7-segment display (display/)
  DisplayController                 8-digit 7-seg multiplexer (display/)
  vio_0                             Xilinx VIO debug core
\end{lstlisting}

\section{Clock domains}

All three clocks are generated from a single MMCM (\texttt{clk\_wiz\_0}) driven by
the 100~MHz board oscillator. Using one MMCM guarantees that all clocks are
phase-coherent and that dynamic phase shifts affect only the selected output.

\begin{longtable}{@{}llll@{}}
  \toprule
  \textbf{Clock} & \textbf{Freq.} & \textbf{Phase} & \textbf{Drives} \\
  \midrule
  \endhead
  \texttt{clk\_sys} & 100~MHz & 0° (fixed) & All RTL logic; FF2 of sensor \\
  \texttt{psclk}    & 100~MHz & 0° + dynamic shift & FF1 of sensor (\texttt{modern\_sensible}) \\
  \texttt{clk\_en}  & 100~MHz & 100° (fixed) & FF3 (alarm latch) of sensor \\
  \bottomrule
\end{longtable}

The \texttt{alarm\_sig} output of \texttt{modern\_sensible} is in the
\texttt{clk\_en} domain. It is passed through a standard 2-FF synchroniser
(\texttt{alarm\_meta / alarm\_sync}) before being used by any
\texttt{clk\_sys} logic.

\section{The aging-sensitive critical path}

The path whose degradation we want to measure is a 16-bit ripple-carry adder
(\texttt{ripple\_adder}) built entirely from LUT primitives:

\begin{lstlisting}[style=verilog, caption={LUT full adder instantiation (lut\_full\_adder.sv)}]
(* DONT_TOUCH = "yes" *)
LUT2_L #(.INIT(4'b0110)) XOR_SUM (.LO(s_wire), .I0(a_i), .I1(b_i));
(* DONT_TOUCH = "yes" *)
LUT3   #(.INIT(8'hE8))   AND_MAJ (.O(c_o),    .I0(a_i), .I1(b_i), .I2(c_i));
\end{lstlisting}

The \texttt{DONT\_TOUCH} attribute is \emph{critical}. Without it, Vivado's
synthesiser replaces the LUT ripple chain with the device's dedicated
\texttt{CARRY4} primitives, which:
\begin{enumerate}
  \item Use a completely different physical path (carry chain routing, not LUT routing).
  \item Are not constrained by the fixed PnR constraints.
  \item Age at a different rate from the LUT fabric.
\end{enumerate}

The result would be that the measured timing degradation does not reflect LUT aging.

\subsection{Fixed placement and routing}

The RTL project uses \emph{fixed placement and routing} (\texttt{fixed\_pnr\_constraints.xdc}).
This file pins every LUT and flip-flop of the sensor to specific BEL locations and
locks the routing. The effect is:

\begin{itemize}
  \item All repeated experiments use exactly the same physical transistors.
  \item There is no run-to-run variation in timing due to the place-and-route algorithm.
  \item Measurements from different experimental sessions are directly comparable.
\end{itemize}

\begin{warningbox}
Never modify the LOC, BEL, or FIXED\_ROUTE constraints in \texttt{fixed\_pnr\_constraints.xdc}
without understanding the consequences. Any change invalidates all previously collected data.
Regenerate from a new checkpoint only after a deliberate PnR run (see Section~\ref{sec:build}).
\end{warningbox}

\section{The metastability sensor (\texttt{modern\_sensible})}

The sensor detects when the critical path output (\texttt{crit\_bit}) is arriving
at a flip-flop too close to the clock edge, i.e., when timing slack approaches zero.

\begin{lstlisting}[style=verilog, caption={modern\_sensible.sv (simplified)}]
// FF1: captures crit_bit on the PHASE-SHIFTED clock (psclk)
FDCE FF1 (.Q(ff1), .C(psclk), .D(in_sensor), ...);

// FF2: captures crit_bit on the REFERENCE clock (sclk = clk_sys)
FDCE FF2 (.Q(ff2), .C(sclk),  .D(in_sensor), ...);

// XOR: detects disagreement between the two captures
LUT2_L #(.INIT(4'b0110)) XOR1 (.LO(xor_out), .I0(ff1), .I1(ff2));

// FF3: latches the XOR result on a third clock (clk_en, 100 deg offset)
FDCE FF3 (.Q(alarm), .C(clk_en), .D(xor_out), ...);
\end{lstlisting}

The operating principle is:

\begin{enumerate}
  \item FF2 always captures \texttt{crit\_bit} reliably on \texttt{clk\_sys}. It represents
        the ``ground truth'' value.
  \item FF1 captures \texttt{crit\_bit} on \texttt{psclk}, which is phase-shifted relative to
        \texttt{clk\_sys} by a controllable amount.
  \item When \texttt{psclk} is close in phase to \texttt{clk\_sys}, FF1 and FF2 see the same
        signal edge and capture the same value $\Rightarrow$ XOR = 0 $\Rightarrow$ no alarm.
  \item As the phase of \texttt{psclk} is shifted (decremented), a point is reached where
        FF1's setup window straddles the transition of \texttt{crit\_bit}. FF1 captures the
        \emph{old} value while FF2 captures the \emph{new} value $\Rightarrow$ XOR = 1
        $\Rightarrow$ alarm fires.
  \item FF3 latches the alarm on \texttt{clk\_en} (100° offset) to avoid hold-time hazards
        in the XOR output.
\end{enumerate}

The number of MMCM phase-shift steps at which the alarm fires is the timing margin measurement.
This is the \textbf{slack} value (\texttt{display\_value}).

\section{Dual-adder design (\texttt{adder\_canary})}

Two structurally identical ripple-adder instances serve different roles:

\subsection{Sensor adder (\texttt{u\_sensor}): timing measurement}

The sensor adder must exercise the full carry chain on every clock cycle to give a
consistent, repeatable timing stimulus. A free-running counter would mostly add small
numbers that do not propagate carry through all 16 stages. Instead, a \texttt{toggle}
flip-flop alternates \texttt{a\_sensor} between two boundary values:

\begin{align}
  \texttt{a\_sensor} &= \begin{cases} 16\texttt{h}5555 & \text{if toggle = 0} \\ 16\texttt{h}5556 & \text{if toggle = 1} \end{cases} \\
  \texttt{b} &= 16\texttt{h}\text{AAAA} \quad(\text{constant})
\end{align}

With these values, $\texttt{sum\_sensor}[15]$ transitions on \emph{every} clock cycle
because the boundary values straddle the carry-out threshold of every stage. This
forces carry propagation through all 16 LUT stages on each cycle, making the timing
stimulus deterministic and maximally stressful on the critical path.

\subsection{Canary adder (\texttt{u\_canary}): functional failure detection}

The canary adder uses a free-running counter input \texttt{a} and the same constant
\texttt{b = 0xAAAA}. Its output is compared against a reference computation
(\texttt{sum\_ref}) implemented in standard synthesis (may use CARRY4). Because
both adders use identical transistors (same placement), the canary detects when
aging has progressed far enough to cause \emph{functional errors}:

\begin{itemize}
  \item \texttt{wrong}: the canary adder output at the moment of the last alarm.
  \item \texttt{correct}: the reference adder output at the same moment.
  \item \texttt{error\_count}: running count of cycles where the two outputs differ (saturates at 0xFFFF).
  \item \texttt{error\_any}: sticky flag, set on the first ever mismatch.
\end{itemize}

\begin{infobox}
The canary metrics are logged to CSV but not displayed in the GUI. They are useful
for post-processing to correlate the onset of timing degradation with eventual
functional failure.
\end{infobox}

\section{The phase-sweep FSM (\texttt{controller\_controller})}

The FSM runs autonomously and continuously. It performs a complete measurement sweep
every few milliseconds without any trigger from the host PC.

\begin{lstlisting}[style=shell, caption={FSM state diagram}]
CHECK_ALARM --> alarm?  --yes--> DONE
     ^                  --no --> INIT_SHIFT
     |                               |
     |              (decrement psclk phase by one step)
     |                               v
     |                          WAIT_SHIFT
     |                               |
     '----(inc_count++)--------------'

DONE --> latch display_value = inc_count
     --> RESET_PHASE

RESET_PHASE --> reset_count >= inc_count? --yes--> IDLE
                                          --no --> WAIT_RESET
                                                       |
              (increment psclk phase, back to 0)       |
              <--(reset_count++)-----------------------'

IDLE --> clear inc_count, reset_count
     --> CHECK_ALARM  (immediately restarts next sweep)
\end{lstlisting}

\textbf{Key insight}: as the device ages, \texttt{crit\_bit} transitions later relative
to the clock edge (path delay increases). The alarm therefore fires after \emph{fewer}
phase-shift steps. Thus \texttt{display\_value} decreases monotonically as the device
degrades. A plot of \texttt{dut\_slack} vs. time showing a downward trend is the
primary evidence of aging.

\subsection{Sticky alarm latch}

Because \texttt{crit\_bit} toggles every clock cycle, the \texttt{alarm} signal from
\texttt{modern\_sensible} is a brief 1--2 cycle pulse rather than a steady level.
\texttt{CHECK\_ALARM} samples \texttt{alarm} for exactly one cycle and would frequently
miss this pulse. A sticky latch captures the first alarm of each sweep and holds it
until IDLE clears it, ensuring the FSM never misses the alarm regardless of pulse
timing.

\section{XADC interface (\texttt{temp\_catcher})}

The Xilinx XADC monitors two internal signals:
\begin{itemize}
  \item \textbf{Temperature} (register \texttt{0x00}): die temperature using an on-chip bandgap diode.
  \item \textbf{VCCINT} (register \texttt{0x01}): internal core supply voltage.
\end{itemize}

The \texttt{temp\_catcher} module polls these registers via the DRP (Dynamic Reconfiguration
Port) interface. A FSM asserts \texttt{DEN} (data enable), waits for \texttt{DRDY}
(data ready), reads \texttt{DO}, then repeats for the next register. A 128-cycle watchdog
retries if \texttt{DRDY} does not arrive (which can happen if DEN is pulsed incorrectly).

\subsection{Conversion formulas}

The 12-bit ADC code (bits [15:4] of the 16-bit register) is converted to physical
units by \texttt{temp\_catcher} in fixed-point arithmetic:

\begin{align}
  T\;[\text{m}^\circ\text{C}] &= \frac{\texttt{ADC\_code}[15:4] \times 503975}{2^{12}} - 273150 \\
  V_{CCINT}\;[\text{mV}] &= \frac{\texttt{ADC\_code}[15:4] \times 3000}{2^{12}}
\end{align}

These values are transmitted in the serial packet as 24-bit integers (millidegrees and millivolts).
The Python application divides by 1000 to recover floating-point Celsius and volts.

\section{UART interface}

\subsection{Trigger (PC \textrightarrow{} FPGA)}

The PC sends the single byte \texttt{0x54} (ASCII `T'). The \texttt{uart\_rx} module
detects this byte and asserts the \texttt{sendin} signal, which causes \texttt{sensor\_stream}
to latch all sensor values and begin transmitting the 15-byte response.

\subsection{15-byte response packet (FPGA \textrightarrow{} PC)}

All fields are Little Endian (LSB first). The UART runs at 9600 baud.

\begin{longtable}{@{}clll@{}}
  \toprule
  \textbf{Byte(s)} & \textbf{Field} & \textbf{Width} & \textbf{Description} \\
  \midrule
  \endhead
  0--2  & \texttt{temp}        & 24-bit & Die temperature (m$^\circ$C), divide by 1000 for $^\circ$C \\
  3--4  & \texttt{sensor}      & 16-bit & Phase-sweep step count (slack) \\
  5--7  & \texttt{vccint}      & 24-bit & VCCINT (mV), divide by 1000 for V \\
  8     & \texttt{failure}     &  8-bit & Bit 0: functional failure latch \\
  9--10 & \texttt{wrong}       & 16-bit & Canary: last wrong adder result \\
  11--12 & \texttt{correct}    & 16-bit & Canary: expected result at alarm time \\
  13--14 & \texttt{error\_count} & 16-bit & Canary: cumulative mismatch count \\
  \bottomrule
\end{longtable}

\begin{warningbox}
The byte layout above is consumed verbatim by \texttt{DUTWorker.poll\_data()} in the
Python application. Never reorder fields or change widths in \texttt{sensor\_stream.sv}
without updating the parser.
\end{warningbox}

\section{Building the bitstream}
\label{sec:build}

\begin{lstlisting}[style=shell, caption={Build commands (from repo root)}]
cd vivado/aging_study_nexys4ddr

# Sanity check (no Vivado needed)
scripts/check_layout.sh

# Create Vivado project (.xpr)
scripts/create_project.sh           # batch mode
scripts/create_project.sh --gui     # open in Vivado GUI

# Build bitstream (creates artifacts/aging_study_nexys4ddr.bit)
scripts/build_bitstream.sh --jobs 8

# Regenerate fixed PnR constraints from a new checkpoint
scripts/build_bitstream.sh --jobs 8 --refresh-ref

# Remove build output
scripts/clean.sh
\end{lstlisting}

The bitstream and timing reports land in \texttt{artifacts/}.
The pre-built bitstream used for App\_2Nexys is in
\texttt{vivado/aging\_study\_nexys4ddr/bitstreams/}.

% =============================================================================
\chapter{Software Architecture}
\label{ch:software}
% =============================================================================

\section{Applications overview}

There are three Python GUI applications. This guide focuses on \textbf{App\_Nexys}
(single DUT, Artix-7 Nexys4 DDR), which is the primary entry point for new students.
The other two are briefly described in Chapter~\ref{ch:otherapps}.

All apps share the same software patterns: a Qt main window, QThread-based workers,
a global configuration module, and a CSV data logger.

\section{Launching the application}

\begin{lstlisting}[style=shell, caption={Running App\_Nexys}]
# Recommended: root launcher (choose at the dialog)
./run.sh

# Direct launch
cd App_Nexys && ./run.sh
\end{lstlisting}

The \texttt{run.sh} script:
\begin{enumerate}
  \item Creates \texttt{.venv/} using the system Python 3 if it does not exist.
  \item Installs dependencies from \texttt{requirements.txt}.
  \item Sets \texttt{LD\_LIBRARY\_PATH} to prefer the venv's Qt libraries over
        system packages (resolves Qt version conflicts on Linux Mint).
  \item Executes \texttt{python App.py}.
\end{enumerate}

\section{Worker architecture (QThread pattern)}

The GUI event loop must remain responsive at all times. All hardware communication
runs in background \texttt{QThread}s. Each worker is a \texttt{QObject} with a
\texttt{QTimer} driving its polling loop. Signals carry data back to the main thread.

\begin{lstlisting}[style=shell, caption={Thread structure}]
Main Thread (Qt event loop)
  |
  +-- ArduinoWorker (QThread)     -- polls Arduino at ~200 ms intervals
  |     QTimer --> poll_data()
  |     Signal: arduino_data_received(dict)
  |
  +-- PSUWorker (QThread)         -- polls IT6502D at ~500 ms intervals
  |     QTimer --> poll_data()
  |     Signal: psu_data_received(dict)
  |
  +-- DUTWorker (QThread)         -- sends 'T', reads 15-byte packet at ~1s
  |     QTimer --> poll_data()
  |     Signal: dut_data_received(dict)
  |
  +-- TestSequencer (QThread)     -- orchestrates, writes CSV, runs safety checks
        QTimer --> log_data_tick()  (1 Hz)
        Receives signals from all three workers
        Runs VCCINT closed-loop trim
        Runs DUT outer temperature loop (every 1800 ticks)
\end{lstlisting}

\section{DUTWorker: reading FPGA data}

On each poll tick, \texttt{DUTWorker}:
\begin{enumerate}
  \item Flushes the serial receive buffer.
  \item Sends the byte \texttt{b'\textbackslash x54'} (0x54 = `T').
  \item Reads exactly 15 bytes.
  \item Unpacks them according to the packet layout (Table in Section~\ref{ch:rtl}).
  \item Converts raw values: \texttt{temp\_raw / 1000.0} $\rightarrow$ \texttt{dut\_temp\_c};
        \texttt{vccint\_raw / 1000.0} $\rightarrow$ \texttt{dut\_volt}.
  \item Emits \texttt{dut\_data\_received} with a dictionary containing all fields.
\end{enumerate}

Boot transients (temp $> 200\,^\circ$C or vccint $> 2.5\,\text{V}$) are
discarded to avoid spurious data while the XADC is settling after power-on.

\section{ArduinoWorker}

The Arduino communicates via ASCII lines over UART at 115200 baud.

\begin{itemize}
  \item \textbf{PC $\rightarrow$ Arduino}: \texttt{SET\_SETPOINT,<value>\textbackslash n}
  \item \textbf{Arduino $\rightarrow$ PC}: \texttt{DATA,<temp>,<sp>,<out>\textbackslash n}
        where \texttt{out} is the PID output (0--100\%).
\end{itemize}

\section{PSUWorker (IT6502D)}

The ITECH IT6502D is accessed via PyVISA with the \texttt{@py} (pyvisa-py) backend.
The VISA resource string must be a USB resource of the form:
\begin{lstlisting}[style=shell]
USB0::0x1AB1::0x0E11::IT6502D300004::INSTR
\end{lstlisting}

To discover the correct string on a given machine:
\begin{lstlisting}[style=shell]
python3 -c "import pyvisa; rm=pyvisa.ResourceManager('@py'); print(rm.list_resources())"
\end{lstlisting}

SCPI commands used:
\begin{itemize}
  \item \texttt{MEAS:VOLT?} --- read output voltage.
  \item \texttt{MEAS:CURR?} --- read output current.
  \item \texttt{VOLT <value>} --- set output voltage.
  \item \texttt{OUTP ON} / \texttt{OUTP OFF} --- enable/disable output.
\end{itemize}

\section{Configuration system}

\texttt{config.py} is a global-mutable-state module. All workers import it directly;
no dependency injection is used. Settings are persisted to \texttt{settings.json}
in the app's working directory.

Key configuration values:

\begin{longtable}{@{}lll@{}}
  \toprule
  \textbf{Parameter} & \textbf{Default} & \textbf{Meaning} \\
  \midrule
  \endhead
  \texttt{DUT\_BAUD}         & 9600          & FPGA UART baud rate \\
  \texttt{ARDUINO\_BAUD}     & 115200        & Arduino UART baud rate \\
  \texttt{PID\_KP}           & 2.78          & Oven PID proportional gain \\
  \texttt{PID\_KI}           & 0.00106       & Oven PID integral gain \\
  \texttt{PID\_KD}           & 5.0           & Oven PID derivative gain \\
  \texttt{VCCINT\_SETPOINT\_V} & 1.0 V       & Target FPGA core voltage \\
  \texttt{VOLTAGE\_KP}       & 0.1 V/V       & VCCINT P-only trim gain \\
  \texttt{MAX\_OVEN\_TEMP\_C} & 130~$^\circ$C & Safety: oven hard limit \\
  \texttt{MAX\_DUT\_TEMP\_C}  & 160~$^\circ$C & Safety: DUT hard limit \\
  \texttt{MAX\_PSU\_CURRENT\_A} & 3.0~A      & Safety: PSU current limit \\
  \bottomrule
\end{longtable}

\section{Data logger}

\texttt{DataLogger} (in \texttt{logger.py}) writes a single CSV file per test run.
Files are named \texttt{<test\_name>\_<YYYYMMDD\_HHMMSS>.csv} and saved to
\texttt{App\_Nexys/test\_logs/}.

Each file begins with a comment block recording the PID parameters, FOPDT model
coefficients, safety limits, and the start timestamp. The data section contains one
row per second with the following columns:

\begin{longtable}{@{}ll@{}}
  \toprule
  \textbf{Column} & \textbf{Description} \\
  \midrule
  \endhead
  \texttt{time\_sec}       & Elapsed seconds since test start \\
  \texttt{oven\_temp\_c}   & Oven air temperature ($^\circ$C) \\
  \texttt{oven\_setpoint\_c} & Current oven setpoint ($^\circ$C) \\
  \texttt{oven\_output\_pct} & PID output to SSR (0--100\%) \\
  \texttt{psu\_voltage\_v} & PSU output voltage (V), measured by PSU \\
  \texttt{psu\_current\_a} & PSU output current (A) \\
  \texttt{dut\_temp\_c}    & FPGA die temperature ($^\circ$C, from XADC) \\
  \texttt{dut\_slack}      & Phase-sweep step count (timing margin) \\
  \texttt{dut\_volt}       & FPGA VCCINT (V, from XADC) \\
  \bottomrule
\end{longtable}

\section{GUI layout}

The main window uses the Catppuccin Mocha dark theme.

\begin{lstlisting}[style=shell, caption={Window layout}]
+--------------------------------------------------------------+
|  [Oven Control + PID info]  [PSU]  [Test Controls]          |  <- top bar
+--------------------------------------------------------------+
|  Tabs:                                                       |
|  [Sensor]  [Temperature]  [Voltage]  [Log]                  |
+--------------------------------------------------------------+
|  (QSplitter)                                                 |
|  Tab content area                                            |
| - - - - - - - - - - - - - - - - - - - - - - - - - - - - -  |
|  Event log (always visible, pinned below tabs)               |
+--------------------------------------------------------------+
|  Status bar: DUT [connected]  Arduino [connected]  PSU [OK]  |
+--------------------------------------------------------------+
\end{lstlisting}

% =============================================================================
\chapter{Communication Protocols}
\label{ch:protocol}
% =============================================================================

\section{DUT protocol (FPGA)}

\begin{itemize}
  \item \textbf{Baud rate}: 9600, 8N1.
  \item \textbf{Trigger}: PC sends single byte \texttt{0x54} (ASCII `T').
  \item \textbf{Response}: FPGA sends 15 bytes, Little Endian, as described in Section~\ref{ch:rtl}.
  \item \textbf{Timing}: response begins within one UART character time (about 1~ms at 9600 baud)
        after the trigger is received.
\end{itemize}

\section{Arduino protocol}

\begin{itemize}
  \item \textbf{Baud rate}: 115200, 8N1.
  \item \textbf{PC $\rightarrow$ Arduino}: \texttt{SET\_SETPOINT,<float>\textbackslash n}
        (e.g. \texttt{SET\_SETPOINT,125.0\textbackslash n}).
  \item \textbf{Arduino $\rightarrow$ PC}: \texttt{DATA,<temp>,<setpoint>,<output>\textbackslash n}
        (e.g. \texttt{DATA,123.45,125.00,47.30\textbackslash n}).
        Sent autonomously at the PID sample interval (5000~ms by default).
\end{itemize}

\section{PSU protocol (IT6502D)}

\begin{itemize}
  \item \textbf{Physical layer}: USB-TMC (USB Test and Measurement Class).
        No baud rate --- full USB speed.
  \item \textbf{Library}: PyVISA with \texttt{@py} backend (pyvisa-py + libusb).
  \item \textbf{Commands}: SCPI (Standard Commands for Programmable Instruments).
\end{itemize}

\begin{warningbox}
The PSU VISA resource string \emph{must} start with \texttt{USB::} or \texttt{USB0::}.
Legacy serial paths (\texttt{/dev/ttyUSB*}) or ASRL strings are rejected at startup by
\texttt{config.load\_config()} to prevent a ``Could not configure port'' error.
\end{warningbox}

% =============================================================================
\chapter{Control Loops}
\label{ch:control}
% =============================================================================

\section{Oven temperature control (PID)}

\subsection{Plant identification}

The oven was modelled as a First-Order Plus Dead Time (FOPDT) system using a
step-test experiment (\texttt{Arduino-ESP/FOPDT\_Step\_Test.ino}):

\begin{equation}
  G(s) = \frac{1.56\,e^{-150.6s}}{1307.2s + 1}
  \label{eq:fopdt}
\end{equation}

Parameters:
\begin{itemize}
  \item $K = 1.56\,^\circ\text{C}/\%$ (steady-state gain)
  \item $\theta = 150.6\,\text{s}$ (dead time)
  \item $\tau = 1307.2\,\text{s}$ (time constant)
  \item $\theta/\tau = 0.115$ (process difficulty ratio)
\end{itemize}

\subsection{SIMC tuning}

PID parameters were computed with the SIMC (Skogestad Internal Model Control)
method, setting the desired closed-loop time constant $\tau_c = \theta$:

\begin{align}
  K_p &= \frac{1}{K} \cdot \frac{\tau}{\tau_c + \theta} = 2.78 \\
  K_i &= \frac{K_p}{\min(\tau,\,4(\tau_c+\theta))} = 0.00106\,\text{s}^{-1} \\
  K_d &= K_p \cdot \frac{\theta}{2} = 5.0\,\text{s}
\end{align}

\begin{warningbox}
Do not change these PID values without conducting a new FOPDT step test.
The parameters are finely tuned to the specific oven and SSR combination in the lab.
Wrong values will cause sustained oscillation or sluggish response.
\end{warningbox}

\subsection{Implementation}

The PID runs on the Arduino at a sample interval of 5000~ms using the ArduPID 1.0.1
library. Key API calls:

\begin{lstlisting}[style=shell, caption={ArduPID 1.0.1 usage}]
pid.setTunings(Kp, Ki, Kd);   // replaces begin() from v0.2.1
pid.setILimits(-100, 100);    // replaces setWindUpLimits()
pid.setDtMs(5000);             // replaces setSampleTime()
output = pid.compute(input);   // takes input; returns output
\end{lstlisting}

The Arduino sends the PID output as a duty cycle percentage to the SSR (Solid-State Relay),
which switches mains power to the oven heating element.

\section{VCCINT closed-loop voltage trim}

The FPGA's internal supply $V_{CCINT}$ is measured by the XADC and included in every
15-byte packet (\texttt{dut\_volt}). The TestSequencer applies a P-only correction
every log tick (1~Hz):

\begin{equation}
  v_{cmd}(k) = v_{cmd}(k-1) + K_{p,V} \cdot \bigl(V_{set} - V_{meas}(k)\bigr)
\end{equation}

with $K_{p,V} = 0.1\,\text{V/V}$. The command is clamped to $[0.0\,\text{V},\;1.5\,\text{V}]$.
This compensates for cable resistance and the PSU's load regulation, ensuring the
FPGA core actually operates at the commanded voltage.

\section{DUT outer temperature loop}

Because the oven air temperature and the FPGA die temperature differ (the die generates
heat internally), a slow outer loop adjusts the oven setpoint to drive the die temperature
to the target. Every 1800 ticks (approximately 30~minutes) the sequencer:

\begin{enumerate}
  \item Computes $e = T_{die,target} - T_{die,measured}$.
  \item If $|e| > 3\,^\circ\text{C}$ (dead-band), shifts the oven setpoint by $\pm 1\,^\circ\text{C}$
        in the direction that reduces the error.
\end{enumerate}

This ensures that the FPGA transistors experience the intended stress temperature
regardless of self-heating variations.

% =============================================================================
\chapter{Setting Up a Fresh Workstation}
\label{ch:setup}
% =============================================================================

\section{System requirements}

\begin{itemize}
  \item OS: Ubuntu 22.04 LTS or Linux Mint 21+ (recommended). Other Debian-based
        distros should work with minor adjustments.
  \item Python: 3.10 or later (system package).
  \item Vivado: 2025.2 or later (required for bitstream builds and App\_2Nexys
        auto-programming; not needed just to run App\_Nexys with a pre-built bitstream).
  \item Disk space: \textasciitilde{}70~GB for Vivado; \textasciitilde{}500~MB for the repo + venv.
\end{itemize}

\section{Step 1 --- Install system dependencies}

\begin{lstlisting}[style=shell, caption={System package installation}]
sudo apt update
sudo apt install -y \
    python3 python3-pip python3-venv \
    libusb-1.0-0 libusb-dev \
    git
\end{lstlisting}

\section{Step 2 --- Clone the repository}

\begin{lstlisting}[style=shell]
git clone https://github.com/your-org/IC-Aging-Environment.git
cd IC-Aging-Environment
\end{lstlisting}

\section{Step 3 --- Install Vivado (if needed)}

Download the Vivado 2025.2 installer from the Xilinx/AMD downloads page.
Run the installer and select at minimum:

\begin{itemize}
  \item \textbf{Vivado} (not just Vitis).
  \item \textbf{Artix-7} device support (under 7 Series).
  \item \textbf{Cable Drivers} (Digilent JTAG over USB).
\end{itemize}

Add Vivado to \texttt{PATH} in \texttt{\textasciitilde/.bashrc}:

\begin{lstlisting}[style=shell]
# Replace the path with your actual Vivado installation directory
echo 'source /opt/Xilinx/Vivado/2025.2/settings64.sh' >> ~/.bashrc
source ~/.bashrc
vivado -version   # should print version info
\end{lstlisting}

\begin{infobox}
If you use App\_2Nexys, you must also update \texttt{VIVADO\_BIN} in
\texttt{App\_2Nexys/config.py} to point to your local Vivado binary. This path
is currently hardcoded to a specific machine in the lab.
\end{infobox}

\section{Step 4 --- Configure udev rules for serial and JTAG devices}

By default, non-root users cannot open USB serial ports or JTAG adapters.
Two steps are needed:

\subsection{Add user to the dialout group}

\begin{lstlisting}[style=shell]
sudo usermod -aG dialout $USER
# Log out and back in for the group change to take effect.
# Verify:
groups | grep dialout
\end{lstlisting}

\subsection{Install Digilent USB cable rules}

The Nexys4 DDR uses an FTDI chip programmed with Digilent's VID/PID.
Vivado ships the required udev rules:

\begin{lstlisting}[style=shell]
# Path may differ; look for install_drivers in your Vivado installation
sudo /opt/Xilinx/Vivado/2025.2/data/xicom/cable_drivers/lin64/install_script/\
  install_drivers/install_drivers
sudo udevadm control --reload-rules
sudo udevadm trigger
\end{lstlisting}

\begin{warningbox}
Without these rules, Vivado cannot program the FPGA via JTAG, and the serial
port may not appear correctly. If you see \texttt{Permission denied} when opening
a serial port, check group membership first.
\end{warningbox}

\section{Step 5 --- Identify serial port assignments}

Each Nexys4 DDR enumerates \emph{two} \texttt{ttyUSB*} ports via its FTDI chip:

\begin{itemize}
  \item \textbf{Lower-numbered port} (e.g., \texttt{ttyUSB0}): JTAG channel. Do not open this.
  \item \textbf{Higher-numbered port} (e.g., \texttt{ttyUSB1}): UART data channel. Use this for the DUT.
\end{itemize}

To identify which port is which:

\begin{lstlisting}[style=shell]
# List USB serial devices with their USB IDs
ls -la /dev/serial/by-id/
# or use udevadm for details
udevadm info -a /dev/ttyUSB1 | grep -E "idVendor|idProduct|serial"
\end{lstlisting}

The Arduino will appear as a separate entry (typically an ACM or CH340 device).

\section{Step 6 --- Identify the PSU VISA resource}

\begin{lstlisting}[style=shell]
python3 -c "
import pyvisa
rm = pyvisa.ResourceManager('@py')
print(rm.list_resources())
"
# Expected output contains something like:
# ('USB0::0x1AB1::0x0E11::IT6502D300004::INSTR',)
\end{lstlisting}

Copy the full USB string; you will paste it into the Setup Dialog.

\section{Step 7 --- First run}

\begin{lstlisting}[style=shell]
# From repo root
./run.sh
\end{lstlisting}

On first launch, \texttt{run.sh} creates the Python virtual environment
(\texttt{App\_Nexys/.venv}) and installs all dependencies from \texttt{requirements.txt}.
This may take a minute. Subsequent launches are instant.

The application opens. If \texttt{settings.json} does not exist, the Setup Dialog
appears automatically.

% =============================================================================
\chapter{Running an Experiment}
\label{ch:running}
% =============================================================================

\section{Pre-experiment checklist}

Before starting a test, verify:

\begin{itemize}
  \item[$\square$] FPGA board is inside the oven with USB cable routed out.
  \item[$\square$] PSU is connected to the FPGA's power header; output is OFF.
  \item[$\square$] Arduino is connected via USB; serial cable to SSR is connected.
  \item[$\square$] SSR wiring to oven is correct (verify polarity and fusing).
  \item[$\square$] NTC thermistor is positioned inside the oven.
  \item[$\square$] All USB devices are enumerated on the PC (\texttt{lsusb}).
\end{itemize}

\section{Setup Dialog}

On first launch (or after clicking \textbf{Setup} in the menu), the Setup Dialog
appears. Fill in:

\begin{longtable}{@{}ll@{}}
  \toprule
  \textbf{Field} & \textbf{Value} \\
  \midrule
  \endhead
  DUT Port & Higher-numbered ttyUSB for the Nexys4 DDR (e.g., \texttt{/dev/ttyUSB1}) \\
  DUT Baud & 9600 (fixed; do not change) \\
  Arduino Port & The Arduino's serial port (e.g., \texttt{/dev/ttyUSB2} or \texttt{/dev/ttyACM0}) \\
  Enable Arduino & Check if the oven is connected \\
  PSU Resource & The USB VISA string discovered in Step~6 \\
  Enable PSU & Check if the IT6502D is connected \\
  VCCINT Setpoint & Target core voltage (V); default 1.0~V; raise to accelerate aging \\
  \bottomrule
\end{longtable}

Click \textbf{Save}. The settings are written to \texttt{settings.json}.

\section{Setting experiment parameters}

In the main window top bar:

\begin{enumerate}
  \item \textbf{Oven setpoint}: type the desired temperature in the oven control group
        (e.g., 125~$^\circ$C). The PID ramp will bring the oven there at 1~$^\circ$C/s.
  \item \textbf{Test name}: enter a descriptive name for the experiment. This becomes
        the CSV filename prefix.
  \item \textbf{Enable Arduino / Enable PSU}: toggle these to match your hardware.
\end{enumerate}

\section{Starting the test}

Click \textbf{Start Test}. The sequencer:
\begin{enumerate}
  \item Enables PSU output and ramps to \texttt{VCCINT\_SETPOINT\_V}.
  \item Sends the oven setpoint to the Arduino.
  \item Begins polling the FPGA at 1~Hz.
  \item Opens a new CSV file with a timestamped name.
\end{enumerate}

The \textbf{Sensor} tab shows the live \texttt{dut\_slack} value. The
\textbf{Temperature} tab shows oven and DUT temperatures.

\section{Monitoring}

Experiments run unattended for days. The app continues logging even if the
display goes to sleep. Recommended practice:

\begin{itemize}
  \item Disable screen lock and display power management on the test machine.
  \item Check in on the CSV daily to verify that new rows are being written.
  \item Monitor \texttt{dut\_temp\_c} in the log: it should be stable within $\pm 1\,^\circ$C
        of the target once the PID has settled (approximately 30~min after start).
\end{itemize}

\section{Safety limits}

The TestSequencer monitors safety limits every tick. If any limit is exceeded,
the test is automatically stopped and the PSU is disabled:

\begin{itemize}
  \item Oven temperature $> 130\,^\circ\text{C}$
  \item DUT die temperature $> 160\,^\circ\text{C}$
  \item PSU current $> 3.0\,\text{A}$
\end{itemize}

An event is logged to the on-screen log panel with a timestamp.

\section{Stopping the test}

Click \textbf{Stop Test}. The sequencer:
\begin{enumerate}
  \item Closes the CSV file and writes an end-of-test footer.
  \item Disables PSU output.
  \item Sends a 0~$^\circ$C setpoint to the Arduino (which effectively stops heating).
\end{enumerate}

Allow the oven to cool before opening it. Do not disconnect USB cables while the
oven is hot.

% =============================================================================
\chapter{Data Analysis}
\label{ch:analysis}
% =============================================================================

\section{CSV file structure}

Log files are in \texttt{App\_Nexys/test\_logs/}. Each file begins with comment
lines (\texttt{\#}) containing metadata, followed by a header row, then data rows
at 1-second intervals.

\section{Loading data in Python (Pandas)}

\begin{lstlisting}[style=python, caption={Loading and plotting dut\_slack over time}]
import pandas as pd
import matplotlib.pyplot as plt

df = pd.read_csv('test_logs/MyTest_20260101_120000.csv', comment='#')

fig, axes = plt.subplots(3, 1, figsize=(12, 8), sharex=True)

axes[0].plot(df['time_sec'] / 3600, df['dut_slack'])
axes[0].set_ylabel('Slack (phase steps)')
axes[0].set_title('Timing Margin vs. Time')

axes[1].plot(df['time_sec'] / 3600, df['dut_temp_c'], label='DUT die')
axes[1].plot(df['time_sec'] / 3600, df['oven_temp_c'], label='Oven air', alpha=0.5)
axes[1].set_ylabel('Temperature (degC)')
axes[1].legend()

axes[2].plot(df['time_sec'] / 3600, df['dut_volt'])
axes[2].set_ylabel('VCCINT (V)')
axes[2].set_xlabel('Time (hours)')

plt.tight_layout()
plt.savefig('aging_result.png', dpi=150)
plt.show()
\end{lstlisting}

\section{Typical result signature}

A successful aging run shows:
\begin{itemize}
  \item \texttt{dut\_slack} decreasing monotonically over the stress period.
  \item \texttt{dut\_temp\_c} stable within $\pm 1\,^\circ$C of the target after the
        initial ramp.
  \item \texttt{dut\_volt} stable within $\pm 5\,\text{mV}$ of the setpoint.
  \item \texttt{error\_count} (in the raw packet; not in the default CSV) remaining at 0
        until functional failure begins.
\end{itemize}

If \texttt{dut\_slack} shows large random jumps, check:
\begin{enumerate}
  \item Temperature stability --- thermal gradients affect timing.
  \item $V_{CCINT}$ stability --- supply noise directly modulates path delay.
  \item USB cable quality --- occasional packet errors produce spurious readings.
\end{enumerate}

\section{Fitting the degradation model}

To extract the aging exponent $n$ from the NBTI power-law:

\begin{lstlisting}[style=python, caption={Fitting slack vs. time to a power-law model}]
import numpy as np
from scipy.optimize import curve_fit

# Use only data after thermal stabilisation (e.g., skip first 3600 s)
mask = df['time_sec'] > 3600
t = df.loc[mask, 'time_sec'].values
s = df.loc[mask, 'dut_slack'].values

# Model: slack(t) = s0 - A * t^n
def model(t, s0, A, n):
    return s0 - A * t**n

p0 = [s.max(), 1.0, 0.15]
popt, _ = curve_fit(model, t, s, p0=p0, maxfev=10000)
s0, A, n = popt
print(f"Initial slack: {s0:.1f}  A: {A:.4f}  n: {n:.4f}")
\end{lstlisting}

The fitted exponent $n$ can be compared against theoretical NBTI predictions
($n \approx 0.1$--$0.25$).

% =============================================================================
\chapter{Other Applications}
\label{ch:otherapps}
% =============================================================================

\section{App\_2Nexys (dual DUT)}

App\_2Nexys operates two Nexys4 DDR boards simultaneously with a shared oven,
two PSUs (IT6502D + Agilent E3634A), and two independent DUT workers. Both
FPGAs are programmed automatically at test start via a Vivado batch session.

Key differences from App\_Nexys:
\begin{itemize}
  \item The DUT outer temperature loop uses the \emph{average} of both DUT temperatures.
  \item Port pairing is done by fixed USB serial IDs (\texttt{USB\_ID\_DUT0}/\texttt{DUT1}).
  \item \texttt{VIVADO\_BIN} is hardcoded in \texttt{App\_2Nexys/config.py} --- update
        this when deploying on a new machine.
  \item The 15-byte DUT protocol is identical to App\_Nexys.
\end{itemize}

\section{App\_FPGAging\_Slack\_Sensor (UltraScale+ / SBCCI target)}

This application targets an UltraScale+ FPGA accessed via an ESP32 UART router
and an STM32 bridge. It supports a different 9-byte packet format from the CROC
FPGA and uses CRC16-Modbus framing for STM32 communication.
This is an advanced configuration used for the SBCCI publication target.

% =============================================================================
\chapter{Troubleshooting}
\label{ch:troubleshooting}
% =============================================================================

\section{Serial port not accessible}

\textbf{Symptom}: \texttt{Permission denied: '/dev/ttyUSB1'} or port not found.

\textbf{Fix}:
\begin{lstlisting}[style=shell]
# Verify group membership
groups | grep dialout
# If dialout is not listed:
sudo usermod -aG dialout $USER
# Then log out and log back in.
\end{lstlisting}

\section{Wrong ttyUSB port selected (JTAG vs UART)}

\textbf{Symptom}: App connects but no data arrives; DUT field stays at zero.

\textbf{Fix}: The lower-numbered port of a Nexys4 DDR is JTAG, not UART.
Always select the \emph{higher-numbered} of the two ports in the Setup Dialog.
Use \texttt{ls -la /dev/serial/by-id/} to inspect the symlinks and identify
port numbers.

\section{App\_2Nexys fails to program FPGAs}

\textbf{Symptom}: Test starts but the FSM logs ``Vivado batch job failed''.

\textbf{Fix}: \texttt{VIVADO\_BIN} in \texttt{App\_2Nexys/config.py} is hardcoded
to an absolute path from another machine. Update it to the local Vivado installation:

\begin{lstlisting}[style=shell]
# Find the vivado binary
which vivado   # if already on PATH
# or
find /opt -name vivado -type f 2>/dev/null | head -5

# Then edit App_2Nexys/config.py:
# VIVADO_BIN = '/opt/Xilinx/Vivado/2025.2/bin/vivado'
\end{lstlisting}

\section{PSU VISA string error}

\textbf{Symptom}: \texttt{config.load\_config()} silently clears the PSU port;
the PSU shows as ``not configured'' even though a port was saved.

\textbf{Cause}: The saved port string does not start with \texttt{USB}.
\texttt{load\_config()} rejects non-USB strings to prevent the IT6502D from
being opened as a serial device.

\textbf{Fix}: Re-run the resource discovery command and paste the full USB string
into the Setup Dialog.

\section{dut\_slack shows erratic jumps}

\textbf{Cause}: Most commonly thermal or voltage instability during the initial warm-up.

\textbf{Fix}:
\begin{enumerate}
  \item Wait at least 30 minutes after pressing Start before drawing conclusions.
  \item Verify \texttt{dut\_temp\_c} has stabilised.
  \item Verify \texttt{dut\_volt} is within 5~mV of the setpoint.
  \item If the XADC is not yet settled after power-on, readings of $> 200\,^\circ$C
        will be silently discarded by \texttt{DUTWorker}.
\end{enumerate}

\section{Vivado timing reports show negative slack on the sensor path}

This is expected and intentional. The sensor path is deliberately placed at its
timing limit. Vivado's timing analysis uses post-P\&R RC delays that may not
exactly match the actual operating condition. The MMCM phase measurement
accounts for the real in-situ timing directly.

% =============================================================================
\appendix
\chapter{FOPDT Step Test Procedure}
% =============================================================================

The FOPDT model of Equation~\ref{eq:fopdt} was identified from a step-test
experiment. To repeat this experiment (e.g., after changing the oven or SSR):

\begin{enumerate}
  \item Flash \texttt{Arduino-ESP/FOPDT\_Step\_Test.ino} to the Arduino.
  \item Open a serial monitor at 115200 baud.
  \item Wait for the oven to reach thermal steady state at a low output (e.g., 20\%).
  \item Apply a step change to a higher output (e.g., 40\%).
  \item Record temperature vs. time until a new steady state is reached.
  \item Fit the FOPDT model to extract $K$, $\theta$, and $\tau$.
  \item Apply SIMC tuning rules to compute new PID gains.
  \item Update \texttt{PID\_KP}, \texttt{PID\_KI}, \texttt{PID\_KD} in
        \texttt{App\_Nexys/config.py}.
\end{enumerate}

% =============================================================================
\chapter{Regenerating Fixed PnR Constraints}
% =============================================================================

The \texttt{fixed\_pnr\_constraints.xdc} file locks all sensor primitives to specific
BEL sites. Regenerate it only after a deliberate new place-and-route run:

\begin{lstlisting}[style=shell]
cd vivado/aging_study_nexys4ddr

# Build a new bitstream AND update the reference checkpoint + constraints
scripts/build_bitstream.sh --jobs 8 --refresh-ref

# Or, if you already have a fresh checkpoint and just want to regenerate XDC:
scripts/extract_fixed_pnr_constraints.sh
\end{lstlisting}

After regeneration, commit both \texttt{fixed\_pnr\_constraints.xdc} and
\texttt{references/fixed\_pnr.dcp}. All previously collected timing data
is now incomparable with future data --- document the PnR change in the
experiment log.

\end{document}
